\section{Decisión de modelo principal y sensibilidades}

A partir de los resultados anteriores, la especificación territorial de trabajo corresponde a un
modelo MNL con \texttt{BIP} como alternativa de referencia para variables comunes, que incorpora
atributos del viaje, variables temporales, demanda y oferta operacional, variables
sociodemográficas del Censo 2024, variables OSM de entorno urbano e infraestructura de acceso, y
proxies históricos de ingreso provenientes de la EOD 2012. Esta estructura permite analizar la
elección entre \texttt{BIP}, \texttt{QR\_OTHER} y \texttt{QR\_RED} combinando atributos
operacionales del viaje con características territoriales del origen.

Dentro de esta estructura se mantienen dos especificaciones candidatas. La primera representa la
dimensión de ingreso EOD mediante proporciones continuas de hogares en los tramos
\texttt{ABC1\_proxy} y \texttt{D+E\_proxy}. La segunda utiliza indicadores de alta concentración
zonal de esos mismos grupos. Ambas especificaciones comparten el resto de variables del modelo y
mantienen patrones sustantivos similares en los bloques de viaje, temporalidad, Censo y OSM.

La decisión entre ambas representaciones no debe basarse únicamente en AIC, BIC o
log-verosimilitud. En modelos con variables territoriales correlacionadas, una variable puede aportar
valor interpretativo aun cuando la mejora incremental de ajuste sea moderada. Por esto, la elección
del modelo de reporte debe combinar tres criterios: estabilidad de signos, coherencia sustantiva e
interpretabilidad para política pública. Bajo ese criterio, la especificación continua permite leer
gradientes territoriales de composición de ingreso, mientras que la especificación dummy permite
identificar zonas con alta concentración relativa de ciertos grupos.

Respecto del bloque OSM, se mantiene la inclusión de refugios de transporte y accesos a Metro. La
variable \path{osm_transport_shelter_yes_density_km2_z} aporta una dimensión de infraestructura de
parada que no está representada por las variables operacionales tradicionales. Por su parte,
\path{osm_railway_subway_entrance_density_km2_z} agrega información sobre accesibilidad física a la
red de Metro, complementando la variable de número de líneas de Metro disponibles en la zona. Aunque
estas variables no deben interpretarse causalmente, sí entregan una lectura territorial relevante
sobre el entorno de acceso al sistema.

Respecto del bloque censal, la especificación conserva variables de composición educacional,
vulnerabilidad habitacional, discapacidad, inmigración, composición por sexo y asistencia a educación
parvularia. Estas variables no describen necesariamente al individuo que realiza el viaje, sino el
contexto territorial de la zona de origen. Su valor está en permitir que el modelo distinga patrones
de adopción de medios de pago asociados a la composición sociodemográfica del territorio.

Como sensibilidad adicional, se evaluó la variable \path{eod2012_numveh_mean_z}, que mide el número
promedio histórico de vehículos por hogar en la zona. Esta variable presentó una señal
estadísticamente informativa, especialmente para \texttt{QR\_OTHER}. Sin embargo, no se incorpora en
las especificaciones candidatas junto con los proxies de ingreso, porque representa una dimensión
conceptualmente muy cercana al nivel socioeconómico y presenta alta colinealidad empírica con las
variables EOD de ingreso. En consecuencia, se considera una sensibilidad útil para diagnóstico, pero
no una variable central del modelo de reporte.

Como sensibilidad computacional y de robustez, también se estimaron modelos equivalentes en Larch con
muestras mayores. En particular, se probaron muestras de 5\% y 10\%, aprovechando que Larch permitió
trabajar con mayor volumen de datos sin los problemas de memoria observados previamente. Los
resultados mantuvieron métricas de ajuste y patrones generales consistentes con la estimación en
Biogeme al 2\%. Esto entrega una confirmación adicional de que las señales principales no dependen
únicamente del submuestreo inicial.

En síntesis, el modelo territorial de trabajo queda definido por una estructura común con variables
de viaje, temporalidad, Censo, OSM y EOD 2012. La decisión metodológica pendiente es seleccionar qué
representación del bloque EOD resulta más adecuada para la tesis: una lectura continua de gradientes
territoriales de ingreso o una lectura discreta basada en zonas de alta concentración relativa. En
ambos casos, la especificación debe entenderse como un modelo principal de trabajo, no como el cierre
definitivo de la investigación.

Aún queda pendiente evaluar otras familias de variables que podrían complementar la lectura
territorial, especialmente aquellas vinculadas con calidad de infraestructura de paraderos y
estaciones, accesibilidad peatonal, seguridad del entorno, disponibilidad de puntos de carga o pago,
y medidas operacionales de confiabilidad, regularidad y ocupación. Estas variables podrían mejorar
la capacidad explicativa del modelo y, sobre todo, aumentar su valor interpretativo para una
autoridad de transporte.
